\subsection*{Broader endpoint search and numerical reproducibility}

The extended study retains the five original ARROW outer folds, the external
220 identities, the strict subset and the 15 frozen response predictions. For each
outer training pool, random inner splits hold out 15\% of each source. The
selection score is the equally weighted mean of public-source and ARROW-source
validation MAE, not an outer-test score. Neural fits use source weights 1 and 3,
respectively. Candidate screening is followed by repeated inner validation,
longer-training challenges up to 1,200 epochs and a separate width challenge.
Each family has a locked configuration and training schedule for each outer
partition. Final fits restart from their specified initialization, restore the
inner-validation rows to training, and average seeds 11, 29, 47, 71 and 101.
Matched no-response fits use the same recipe and schedule. The original tree and
frozen-encoder/tree controls use seeds 11, 29 and 47; common-three-seed comparisons
are also reported. Supplementary Methods~S18 specifies the search and its limits.

MoLFormer representations are generated with the encoder's deterministic evaluation
option in addition to evaluation mode. An audit found that evaluation mode alone
left stochastic attention projections active in the earlier cache. We regenerated
all study and stress representations from the same checkpoint and repeated both
the neural-head and fixed-tree comparisons. Only these corrected MoLFormer results
enter the manuscript's comparisons. Reload, repeated-query and reordered-batch
checks verify prediction replay; the original SolvAI--ExtraTrees release does not
use MoLFormer and is unaffected. The discrepancy and corrected hashes are recorded
in Supplementary Methods~S19 and Data~6.

TabPFN-3.5 uses the pinned full and Fast checkpoints with \texttt{tabpfn} 9.1.0,
run locally under the provider's non-commercial license. We do not fine-tune the
foundation model. Eight candidates combine two checkpoints with four feature
recipes; training-side challenges compare eight versus 32 ensemble passes and
single versus repeated ARROW context rows. Context repetition represents source
emphasis but is not mathematically equivalent to weighted neural loss. Splits
precede repetition. The pretrained model's native preprocessing, including
feature selection, is retained; this is not an entirely tree-free preprocessing
pipeline. Fitted contexts contain training labels and are not redistributed.

\subsection*{Size, series and predictive-uncertainty diagnostics}

The size diagnostic holds out the largest 15\% within each source, with stable
identity ordering for ties. Every model uses the same 1,160/205 split and three
seeds; no held-out label is used in that refit. Configurations come from the locked
external-fit recipe. Because the held-out labels were accessible in earlier
random-fold development, this is an exploratory conditional test of size transfer.
The source models and their source-level out-of-fold predictions are not rebuilt
around this size split; exclusion from the endpoint refit does not imply exclusion
from all source supervision (Supplementary Methods~S5 and S20).
The unlabelled grid comprises 30 linear alkanes and 13 members each of capped glycine
and alanine series, including the common zero-residue molecule. Series identities,
source overlap, fingerprint counts and all predictions are released.

For uncertainty evaluation, we transform each TabPFN distribution to the common
raw-target scale and form an equally weighted mixture across five seeds. Central
80\% and 95\% intervals are quantiles of this mixture, not averages of seedwise
quantiles. Empirical coverage is the fraction of external labels inside each
interval. We also report widths and verify the mixture mean against stored point
predictions. The external cohort is used only to evaluate these intervals, not
to calibrate them. Supplementary Methods~S20--S21 details these diagnostics.
